\scanpage{144}
\mspage{V.16}
% archive sheet number: 143

\deleted{3.4.}\inserted{4.} \deleted{Les cinq} \uline{Situation locale des cinq schémas.}

4.1 Soient \deleted{$\unsure{\bar x\longmapsto x}$} $x\in X_0$, $\bar x$ un pt géométrique sur $x$, $X'$ le localisé strict de $X$ en $\bar x$, posons
\[
\tag*{(4.1.1)}
V=X'-\{\bar x\},\quad V_s=V\times_S\{s\},\quad
V_{\unsure{\eta|y}}=V\times_S\eta=V-V_s,\quad
V_{\bar\eta}=V\times_S\bar\eta=V_\eta\times_\eta\bar\eta.
\]
En plus de ces quatre schémas remarquables, nous en introduisons un cinquième. \deleted{à} \inserted{Pour \unsure{cela}, considérons} \deleted{\unsure{le schéma}} le normalisé $\bar S$ de $S$ dans $\bar K$ :
\[
\tag*{(4.1.2)}
\bar S\longrightarrow S,\qquad \bar S=\varprojlim S_\alpha
\]
où les $S_\alpha$ sont les normalisés de $S$ dans les extensions finies \deleted{\unsure{du point}} $K_\alpha$ de $\bar K/K$. \deleted{+} On posera
\[
\tag*{(4.1.3)}
\bar V=V\times_S\bar S=\varprojlim V_\alpha,\qquad V_\alpha=V\times_S S_\alpha.
\]
\deleted{\unsure{Soit alors} $\bar V_s=\unsure{V_s'}$.}
Evidemment on a un isomorphisme
\[
\deleted{V_{\bar\eta}\overset{\unsure{\mathrm{dfn}|\mathrm{df}}}{=}\bar V\times_S\eta
\simeq\bar V_\eta\overset{\mathrm{dfn}}{=}\bar V\times_{\bar S}\bar\eta\simeq V_{\bar\eta}}
\]
\margin{\diagram{\begin{tikzcd}[ampersand replacement=\&,column sep=small,row sep=small]
\& V\arrow[ddd] \& \& \bar V\arrow[ll]\arrow[ddd] \& \\
\& \& V_\eta\arrow[ul]\arrow[rr]\arrow[ddd] \& \& V_{\bar\eta}\arrow[ul]\arrow[ddd] \\
s\arrow[dr] \& \& \& \bar s\arrow[d] \& \\
\& S \& \& \bar S\arrow[ll] \& \\
\& \& \eta\arrow[ul] \& \& \bar\eta\arrow[ll]\arrow[ul]
\end{tikzcd}}}
\margin{\diagram{\begin{tikzcd}[ampersand replacement=\&,column sep=small,row sep=small]
\deleted{V}\arrow[r]\arrow[d] \& \deleted{\bar V}\arrow[d] \\
\deleted{S}\arrow[r] \& \deleted{\bar S}
\end{tikzcd}}}
\deleted{$V_{\bar\eta}$}
\[
\tag*{(4.1.4)}
V_{\bar\eta}\lhook\joinrel\longrightarrow\bar V,
\]
\unsure{identifiant|identifie} $V_{\bar\eta}$ à la fibre générique de $\bar V$ sur $S$ [resp. sur $\bar S$], et le complémentaire \deleted{$\unsure{V_\eta}$} $\bar V-V_{\bar\eta}$ \deleted{\unsure{est aussi} $\unsure{V^{\prime\prime}}$} \inserted{s'identifie à} \deleted{$V_s\unsure{\times_S S}$}
\[
\tag*{(4.1.5).}
V_{\bar s}=V_s\otimes_{k(s)}k(\bar s)
\]
où $\bar s$ est le point fermé de $\bar S$. Notons que la projection
\[
\tag*{(4.1.6)}
V_{\bar s}\longrightarrow V_s
\]
induit une \uline{équivalence} des topos étales \inserted{par} SGA4 \uline{VIII}. (car $k(s)$ est sép. clos et $k(\bar s)$ en est une extension radicielle), et on \deleted{\unsure{fera}} \inserted{pourra} donc identifier $V_s$ et $V_{\bar s}$.
